\subsection{Variational energies and feasible spaces}
\label{sec:models}

The three \(p\)-energies are defined in \cref{eq:pure-energy,eq:reg-energy,eq:reaction-energy}. The
remaining models have the following energies, written for \(d\leq2\):
\begin{align}
 \E_{\rm lin}(u)&=\int_\Omega\left(\frac12\abs{\nabla u}^2+
 \frac12u^2\right)\dd x-\ip fu,
 \label{eq:linear-energy}\\
 \E_{\rm cub}(u)&=\int_\Omega\left(\frac12\abs{\nabla u}^2+
 \frac14u^4\right)\dd x-\ip fu,
 \label{eq:cubic-energy}\\
 \E_{\sinh}(u)&=\int_\Omega\left(\frac12\abs{\nabla u}^2+
 \cosh u-1\right)\dd x-\ip fu.
 \label{eq:sinh-energy}
\end{align}
The pairing is the energy-space duality; for a regular load it is the corresponding integral.
The additive \(-1\) in \cref{eq:sinh-energy} normalizes the
potential and does not change the Euler--Lagrange equation.

\begin{assumption}[Data and feasible spaces for the six models]
\label[assumption]{ass:model-data}
The following model-specific conditions are imposed.
\begin{enumerate}[label=\textup{(\roman*)}]
 \item For the linear and cubic models, \(f\in (H^1(\Omega))^*\). For the sinh model, the convex
 integral functional has a load \(f\in(H^1(\Omega))^*\), and the
 integral functional is allowed to take the value \(+\infty\), its effective domain is nonempty, and
 local derivative/Hessian statements are restricted to an \(L^\infty\)-bounded neighborhood.
 \item For the pure and regularized \(p\)-energies, \(f\in X_{p,\circ}^*\), or equivalently the Neumann
 load annihilates constants before restriction to \(X_{p,\circ}\).
 \item For the reaction \(p\)-energy, \(f\in(W^{1,p})^*\). In the regularized model, the parameter
 \(\eps\) in \cref{eq:reg-energy} is fixed and positive.
\end{enumerate}
\end{assumption}

\begin{proposition}[Well-posedness]
\label{prop:wellposedness}
Under \Cref{ass:model-data}, \cref{eq:linear-energy,eq:cubic-energy,eq:sinh-energy} have unique minimizers
in \(H^1(\Omega)\); \cref{eq:pure-energy,eq:reg-energy} have unique minimizers in \(X_{p,\circ}\); and
\cref{eq:reaction-energy} has a unique minimizer in \(W^{1,p}(\Omega)\).
\end{proposition}

\begin{proof}
The positive zero-order term in \(\E_{\rm lin}\) controls the complete \(H^1\)-norm. For the cubic
energy, the \(L^4\) term controls constants, while its combination with the gradient and Young's
inequality gives coercivity. Strict convexity follows from the gradient term and \(t\mapsto t^4\). For
the sinh energy, \(\cosh t-1\geq t^2/2\) gives coercivity. Along a weakly convergent \(H^1\) sequence,
compactness in \(L^2\), almost-everywhere convergence of a subsequence, and Fatou's lemma give lower
semicontinuity of the potential.

For the pure and regularized \(p\)-energies, the zero-mean Poincar\'e inequality turns the gradient
seminorm into a norm, and the diffusion densities are strictly convex. For the reaction model, the
gradient and scalar \(p\)-terms control the full \(W^{1,p}\)-norm. The direct method gives existence in
each case, and strict convexity gives uniqueness.
\end{proof}

The local Hessian structure of each model is needed for the transfer estimates of
\Cref{sec:c5-hessian}. The bilinear forms at the exact solution are
\begin{align*}
 a_{\rm lin}(v,w)&=\int_\Omega(\nabla v\cdot\nabla w+vw)\dd x,\\
 a_{\rm cub,*}(v,w)&=\int_\Omega(\nabla v\cdot\nabla w+3(u^*)^2vw)\dd x,\\
 a_{\sinh,*}(v,w)&=\int_\Omega(\nabla v\cdot\nabla w+\cosh(u^*)vw)\dd x.
\end{align*}
For a \(p\)-diffusion density, the Hessian matrix in gradient variables is
\[
 \abs\xi^{p-2}I+(p-2)\abs\xi^{p-4}\xi\otimes\xi
\]
at nonzero \(\xi\) in the pure case, with \(\abs\xi^2\) replaced by
\(\eps^2+\abs\xi^2\) in the regularized case. The local Hessian statements that follow concern
\(p\geq2\).

\begin{proposition}[Sufficient continuous local Hessian conditions]
\label{prop:model-hessian}
In dimensions \(d\leq2\), the cubic Hessian is locally Lipschitz as a bilinear form on \(H^1\), and it
is coercive at any solution \(u^*\not\equiv0\). The sinh Hessian is coercive and is locally Lipschitz on
an \(L^\infty\)-bounded neighborhood, with the state perturbation measured in \(H^1\). For fixed
\(\eps>0\), the regularized \(p\)-Hessian is uniformly elliptic, with lower bound proportional to
\(\eps^{p-2}\); its local Lipschitz continuity is valid in the bounded-gradient topology
\(W^{1,\infty}\), not in the \(H^1\) topology asserted for the cubic model.
\end{proposition}

\begin{proof}
For the cubic multiplier,
\[
 \left|\int_\Omega3(u^2-(u^*)^2)vw\dd x\right|
 \leq C\norm{u-u^*}_{L^4}\norm{u+u^*}_{L^4}
 \norm v_{L^4}\norm w_{L^4},
\]
and \(H^1\hookrightarrow L^4\) for \(d\leq2\). If \(u^*\not\equiv0\), coercivity follows by
contradiction: a unit \(H^1\)-sequence with vanishing quadratic form has gradients tending to zero and,
by compact embedding, converges in \(L^4\) to a constant; the weighted term then forces that constant
to vanish. On \(\abs u,\abs{u^*}\leq R_\infty\), the mean
value theorem gives
\(\abs{\cosh u-\cosh u^*}\leq\sinh(R_\infty)\abs{u-u^*}\), after which the chosen local
multiplication estimate applies. In the regularized model, the radial and tangential eigenvalues are
comparable to \((\eps^2+\abs\xi^2)^{(p-2)/2}\); the lower bound is therefore proportional to
\(\eps^{p-2}\). The matrix-valued mean value theorem gives Hessian stability when the gradients and
their perturbation are controlled in \(L^\infty\). Such control is not a consequence of a bounded
\(H^1\) neighborhood.
\end{proof}

\begin{proposition}[Restricted stability for the regularized \(p\)-energy]
\label{prop:regularized-restricted}
Let \(Y_M\subset W^{1,\infty}(\Omega)\) be finite dimensional and let \(\norm\cdot_{0,M}\) be any norm
on \(Y_M\). For fixed \(\eps>0\), bounded subsets of \(Y_M\) have finite restricted constants
\(L_M\) on the corresponding finite-dimensional segments. If the reference norm is chosen so that the gradient seminorm is a
norm on the feasible space, then \(\mu_M>0\). The resulting constants may depend on \(M\), the inverse
inequality on \(Y_M\), and \(\eps\).
\end{proposition}

\begin{proof}
Norm equivalence on \(Y_M\) converts a \(\norm\cdot_{0,M}\)-bounded set into a
\(W^{1,\infty}\)-bounded set. The bounded-gradient Hessian estimate in the preceding proof then gives
finite \(L_M\). Uniform ellipticity for fixed \(\eps\), together with the assumed treatment of constant
directions, gives \(\mu_M>0\). No constant uniform in \(M\) or \(\eps\downarrow0\) follows.
\end{proof}

\begin{table}[tbp]
\centering
\small
\caption{Hessian structure and the conclusion available for each experimental energy.}
\label{tab:model-classification}
\begin{tabularx}{\textwidth}{@{}>{\raggedright\arraybackslash}p{0.16\textwidth}>{\raggedright\arraybackslash}p{0.31\textwidth}>{\raggedright\arraybackslash}X@{}}
\toprule
Model & Frozen Hessian & Scope of the conclusion \\
\midrule
Linear & Constant \(H^1\) inner product & CGA and energy-inner-product OGA coincide under exact correction, the same dictionary, and the same selection rule. \\
Cubic & Gradient form plus \(3(u^*)^2\) & Locally stable for \(d\leq2\) and coercive when \(u^*\not\equiv0\); transfer still needs relative score resolution. \\
Sinh & Gradient form plus \(\cosh(u^*)\) & Coercive; Hessian Lipschitz is stated on an \(L^\infty\)-bounded neighborhood, and the numerical use is finite-dimensional. \\
Regularized \(p\) & Uniformly elliptic for fixed \(\eps>0\) & Continuous transfer needs bounded-gradient stability; finite-dimensional transfer has explicit \(M\)-dependent constants. \\
Pure \(p\), \(p>2\) & Weight can vanish at \(\nabla u^*=0\) & Use global \(p\)-geometry or the explicit finite-range estimate. \\
Reaction \(p\), \(p>2\) & Gradient and scalar weights can both vanish & The reaction restores full-space energy coercivity, not uniform frozen-Hessian coercivity. \\
\bottomrule
\end{tabularx}
\end{table}

Energy coercivity and frozen-Hessian coercivity are distinct. A positive reaction term does
not imply that the frozen \(p\)-Hessian is uniformly positive at points where both the solution and its
gradient vanish, and the continuous degeneracy of the pure model does not preclude the restricted
coercivity in the corresponding finite-dimensional metric.
